\documentclass[11pt]{article}
\usepackage[margin=1in]{geometry}
\usepackage{mathpazo}
\usepackage{titlesec}
\usepackage{enumitem}
\usepackage{hyperref}

\hypersetup{colorlinks=true, linkcolor=blue, urlcolor=blue}
\titlespacing*{\section}{0pt}{12pt}{6pt}
\titlespacing*{\subsection}{0pt}{8pt}{4pt}
\setlist{nosep, leftmargin=2em}

\title{\textbf{Supplementary Material:\\[4pt] \large SVA Suite Rationale and Categorization}}
\author{Achyuta Ganesh}
\date{}

\begin{document}
\maketitle

\noindent This supplementary document details the final 30 SystemVerilog Assertions (SVAs) chosen for the Gold-Standard verification suite. Per the evaluation rubric, each property explicitly details what it protects, what bug would slip past without it, and how a violation manifests.

\section*{A) Local State: Reset}
\begin{itemize}
    \item \textbf{P20a (Reset reads zero):} 
    \begin{itemize}
        \item Ensures the pipeline is cleanly initialized at the clock edge.
        \item Without it, a dirty reset condition could leave residual values in the pipeline buffers.
        \item As a result upon boot, CPU would execute garbage instructions.
    \end{itemize}
    
    \item \textbf{P20b (Asynchronous clear):} 
    \begin{itemize}
        \item Ensures the CPU state clears immediately upon reset, independent of the clock.
        \item Without it, a synchronous-only reset could allow a corrupted state to persist (undetected) till the clock edge atleast.
    \end{itemize}
\end{itemize}

\section*{B) Load-Use Hazard $\rightarrow$ Stall}
\begin{itemize}
    \item \textbf{P1 (Happens when it must):} 
    \begin{itemize}
        \item Protects data dependencies between memory loads and immediate ALU uses.
        \item Without it, the hazard detection unit could fail to recognize a dependency and let the pipeline continue.
        \item This would manifest as the dependent instruction computing using stale register data instead of the newly loaded memory data.
    \end{itemize}

    \item \textbf{P1b (Freezes PC \& IF/ID):} 
    \begin{itemize}
        \item Prevents instruction fetch from advancing during a stall.
        \item Without it, a broken stall signal could let the PC increment while the rest of the pipeline waits.
        \item The CPU would skip instructions entirely.
    \end{itemize}

    \item \textbf{P2 (ID/EX clears to bubble):} 
    \begin{itemize}
        \item Injects a NOP to separate the stalled instruction from the executing load.
        \item Without it, a stall might freeze the front half of the pipeline but fail to clear the ID/EX control signals.
        \item The instruction currently sitting in ID/EX would execute twice, corrupting architectural state.
    \end{itemize}

    \item \textbf{P5 (Stalls exactly 1 cycle):} 
    \begin{itemize}
        \item Ensures pipeline efficiency by strictly bounding the load-use penalty.
        \item Without it, a sticky hazard flag could cause a multi-cycle stall.
        \item The CPU would freeze unnecessarily for multiple cycles on a single load-use hazard.
    \end{itemize}
\end{itemize}

\section*{C) Pipeline Handoffs}
\begin{itemize}
    \item \textbf{P18 (PC updates correctly):} 
    \begin{itemize}
        \item Protects the fundamental sequential execution of the program.
        \item Without it, the PC could skip or repeat addresses during normal operation.
    \end{itemize}

    \item \textbf{P3 (IF/ID frozen during stall):} 
    \begin{itemize}
        \item Makes sure the contents of the IF/ID remain forzen during stall
        \item The currently stalled instruction would be overwritten and permanently lost.
        \item P1b checks the control wires to make sure IF/ID is frozen but this checks the actual reg data against the previous clock cycle to gurantee it froze perfectly.
    \end{itemize}

    \item \textbf{P28 (IF/ID captures if no stall):} 
    \begin{itemize}
        \item Protects the normal propagation of fetched instructions.
        \item Without it, a stuck stall flag could prevent new instructions from entering the pipeline.
        \item While P3 checks that during a stall the regs are frozen this makes sure it captures the new data correctly if there is no stall.
    \end{itemize}

    \item \textbf{P25 (Data enters ID/EX if no hazard):} 
    \begin{itemize}
        \item Protects normal instruction flow from Decode to Execute.
        \item Similar to previous, P2 made sure a bubble (NOP) is inserted un the case of stall and this gurantees its working if there is no stall
    \end{itemize}

    % \item \textbf{P34a (MEM to WB writes once):} 
    % \begin{itemize}
    %     \item Protects the processor from committing the same instruction multiple times
    %     \item Without it, a stalling MEM stage could cause the WB stage to commit the same data multiple times.
    %     \item The \texttt{wb\_done} signal would assert repeatedly for a single instruction, corrupting the commit trace.
    % \end{itemize}
\end{itemize}

\section*{D) Backpressure}
\begin{itemize}
    \item \textbf{P42 (Req maintains value):} 
    \begin{itemize}
        \item Makes sure req maintains its valye until ready
        \item Without it, the CPU could change the address or data midway through a wait state.
        % \item The memory controller would receive a fragmented or corrupted write request.
    \end{itemize}

    \item \textbf{P43 (Mem stall freezes pipeline):} 
    \begin{itemize}
        \item Makes sure that a Mem stall correctly freezes the pipeline
        \item Without it, IF, ID, and EX could advance while MEM is blocked.
        \item Instructions currently in EX would overwrite the instruction stuck in MEM, destroying it.
    \end{itemize}

    \item \textbf{P44 (Pipeline not stalled forever):} 
    \begin{itemize}
        \item Protects the liveness of the processor against deadlocks.
        % \item Without it, a broken memory controller or logic loop could hold the wait state indefinitely.
        \item The \texttt{ready} signal would never assert, holding the wait state indefinitely.
    \end{itemize}
\end{itemize}

\section*{E) Forwarding}
\begin{itemize}
    \item \textbf{P6 (Youngest older producer wins):} 
    \begin{itemize}
        \item Protects data dependency priorities when multiple stages hold valid data for the same register.
        \item Without it, a double-data hazard could incorrectly prioritize MEM/WB over the newer EX/MEM data.
        \item The ALU would compute using older, stale data instead of the most recent result.
    \end{itemize}

    \item \textbf{P38 (Live operand data correctness):} 
    \begin{itemize}
        \item Ensures that forwarding happens when it must and the data entering the ALU during forwarding = what the reg file would hold if we did not have the pipeline
        \item Without it, a wiring bug could slip past where the multiplexer makes the correct decision, but pulls data from the wrong wire
        % \item The property would fail because the ALU inputs diverge from the ideal, unpipelined register state.
    \end{itemize}
\end{itemize}

\section*{F) Branch, Commits, and Memory}
\begin{itemize}
    \item \textbf{P8 (Flushing happens when it must):} 
    \begin{itemize}
        \item Protects the control flow logic from executing mispredicted instructions.
        % \item Without it, the branch evaluator could detect a jump but fail to trigger the flush wire.
        \item The two wrong-path instructions in IF and ID would incorrectly execute and commit.
    \end{itemize}

    \item \textbf{P8b (Flushed instr has no side effects):} 
    \begin{itemize}
        \item Protects the architectural state (reg/mem writes) from changing due to a flushed instr.
        \item Without it, a flush could zero the instruction but leave \texttt{RegWrite} or \texttt{MemWrite} high.
        \item A squashed instruction would still alter a register or write garbage to memory.
    \end{itemize}

    % \item \textbf{P34b (Commit happens exactly once):} 
    % \begin{itemize}
    %     \item Protects the instruction retirement cadence under backpressure.
    %     \item Without it, a stalled memory instruction could trigger multiple commits while waiting for \texttt{ready}.
    %     \item The scoreboard would flag a mismatch as the CPU commits the same instruction over multiple clock cycles.
    % \end{itemize}

    \item \textbf{P35 (Committed value comes from right source):} 
    \begin{itemize}
        \item Makes sure the final commit value always comes from the right source (ALU/mem)
        \item Without it, a \texttt{lw} instruction could accidentally write ALU results, or an \texttt{add} could write memory data.
        % \item The register file would update with fundamentally incorrect computation outputs.
    \end{itemize}

    % \item \textbf{P36 (Branch evaluated using forwarded data):} 
    % \begin{itemize}
    %     \item Protects data dependencies involving \texttt{JZ} branch evaluations.
    %     \item Without it, the branch comparator could read stale register file data instead of forwarded EX/MEM data.
    %     \item The CPU would branch (or fail to branch) based on outdated conditions, breaking control flow.
    % \end{itemize}

    \item \textbf{P39a (Right \$rs data to memory):} 
    \begin{itemize}
        \item Makes sure the right forwarded data is written to memeory
        \item P38 protects forwarding to the ALU not memory.
        \item A \texttt{sw} instruction would write stale or incorrect data into the memory.
    \end{itemize}

    \item \textbf{P39b (Store lands at right address):} 
    \begin{itemize}
        \item Protects memory addressing calculation.
        \item Without it, an ALU bug could calculate the wrong offset for the memory write.
        \item The CPU would overwrite unrelated data structures in RAM.
    \end{itemize}

    \item \textbf{P39c (No writes to RAM during stall):} 
    \begin{itemize}
        \item Protects data memory from writes while the pipeline waits.
        \item Without it, \texttt{MemWrite} signal could stay turned on while the target memory address could randomly change during the freeze.
        \item The processor would write invalid data across random memory addresses while frozen.
    \end{itemize}
\end{itemize}

\section*{H) Environment \& Architectural Invariants}
\begin{itemize}
    \item \textbf{P51 (\$zero is immutable):} 
    \begin{itemize}
        \item Protects the MIPS architectural guarantee that register 0 cannot be altered.
        \item Without it, a software instruction like \texttt{add \$0, \$1, \$2} could successfully change the hardware register.
        \item Subsequent instructions reading \texttt{\$zero} would receive non-zero values.
    \end{itemize}

    \item \textbf{P40 (Write-through bypass):} 
    \begin{itemize}
        \item Protects against read-after-write hazards in the same clock cycle within the register file.
        \item It is basically another "forwarding"
        \item The ID stage would read stale data, requiring an unnecessary pipeline stall.
    \end{itemize}

    \item \textbf{P26 (Reading \$zero returns 0):} 
    \begin{itemize}
        \item Protects the output wires of the register file.
        \item Without it, the register file could store 0 but physically output noise on the \texttt{\$zero} read port.
        \item The ALU would receive garbage bits when attempting to use \texttt{\$zero}.
    \end{itemize}

    \item \textbf{P12 (PC is multiple of 4):} 
    \begin{itemize}
        \item Protects instruction alignment requirements.
        \item Without it, a branch could calculate an unaligned address (e.g., \texttt{PC = 0x03}).
        \item The instruction memory would return misaligned, corrupted instruction bits.
    \end{itemize}

    \item \textbf{A9 (PC bounded to IMEM):} 
    \begin{itemize}
        \item Protects the formal solver from testing impossible physical limits.
        \item Without it, the solver could inject out-of-bounds addresses (this is an assumption, not an assertion).
        \item The solver would falsely report a hardware bug when the PC exceeds physical memory bounds.
    \end{itemize}

    \item \textbf{P17 (Data accesses bounded to DMEM):} 
    \begin{itemize}
        \item Protects the memory array from out-of-bounds requests.
        \item Without it, an ALU offset could cause a read/write beyond the top of memory.
        \item The simulation would crash or return \texttt{X} (undefined) states.
    \end{itemize}
\end{itemize}

\end{document}